\documentclass{article}
\usepackage{amsmath,amssymb,amsthm}
\usepackage{algorithm,algorithmic}
\usepackage{bm}
\usepackage{hyperref}
\usepackage{enumitem}
\newtheorem{theorem}{Theorem}
\newtheorem{lemma}{Lemma}
\newtheorem{assumption}{Assumption}
\newcommand{\E}{\mathbb{E}}
\newcommand{\R}{\mathbb{R}}
\newcommand{\norm}[1]{\left\|#1\right\|}
\begin{document}

\section*{Algorithm Name}
\textbf{Clipped Gradient Descent with Momentum (CGDM)}  

The method applies element–wise $\ell_2$‑norm clipping to the stochastic gradient, then updates the parameters using a classical momentum accumulator.

\section*{Mathematical Formulation}
Let $f:\R^d\rightarrow\R$ be the (possibly non‑convex) objective and let $\{z_t\}_{t\ge 1}$ denote an i.i.d. data stream.  
At iteration $t$ we compute a stochastic gradient
\[
g_t^{\text{raw}} \;:=\; \nabla f(x_t;z_t)\in\R^d .
\]

\paragraph{Gradient clipping}
\[
\boxed{\;
g_t \;:=\; \frac{g_t^{\text{raw}}}{\max\!\Bigl(1,\;\frac{\norm{g_t^{\text{raw}}}}{C}\Bigr)}\;},
\qquad C>0\;\text{(clipping threshold)}.
\]
Thus $\norm{g_t}\le C$ for all $t$.

\paragraph{Momentum update}
\[
\begin{aligned}
m_t &:= \beta\, m_{t-1} + (1-\beta)\, g_t, \qquad m_0:=0, \\
x_{t+1} &:= x_t - \eta\, m_t ,
\end{aligned}
\]
where $\eta>0$ is the stepsize and $\beta\in[0,1)$ the momentum coefficient.

\section*{Pseudocode}
\begin{algorithm}[H]
\caption{Clipped Gradient Descent with Momentum (CGDM)}
\begin{algorithmic}[1]
\STATE \textbf{Input:} stepsize $\eta>0$, momentum $\beta\in[0,1)$, clipping threshold $C>0$,\\
\ \ \ \ \ \ \ \ \ \ \ \ \ \ \ \ \ \ \ \ \ \ \ \ \ \ \ \ \ \ \ \ \ \ \ \ \ \ \ \ \ \ \ \ \ \ \ \ \ \ \ \ \ \ \ \ \ \ \ \ \ \ \ \ \ \ \ \ \ \ \ \ \ \ \ \ \ \ \ \ \ \ \ \ \ \ \ \ \ \ \ \ \ \ \ \ \ \ \ \ \ \ \ \ \ \ \ \ \ \ \ \ \ \ \ \ \ \ \ \ \ \ \ \ \ \ \ \ \ \ \ \ \ \ \ \ \ \ \ \ \ \ \ \ \ \ \ \ \ \ \ \ \ \ \ \ \ \ \ \ \ \ \ \ \ \ \ \ \ \ \ \ \ \ \ \ \ \ \ \ \ \ \ \ \ \ \ \ \ \ \ \ \ \ \ \ \ \ \ \ \ \ \ \ \ \ \ \ \ \ \ \ \ \ \ \ \ \ \ \ \ \ \ \ \ \ \ \ \ \ \ \ \ \ \ \ \ \ \ \ \ \ \ \ \ \ \ \ \ \ \ \ \ \ \ \ \ \ \ \ \ \ \ \ \ \ \ \ \ \ \ \ \ \ \ \ \ \ \ \ \ \ \ \ \ \ \ \ \ \ \ \ \ \ \ \ \ \ \ \ \ \ \ \ \ \ \ \ \ \ \ \ \ \ \ \ \ \ \ \ \ \ \ \ \ \ \ \ \ \ \ \ \ \ \ \ \ \ \ \ \ \ \ \ \ \ \ \ \ \ \ \ \ \ 
\STATE Initialize $x_1\in\R^d$, $m_0\gets 0$.
\FOR{$t=1,2,\dots,T$}
    \STATE Sample $z_t$ and compute raw stochastic gradient $g_t^{\text{raw}}\gets\nabla f(x_t;z_t)$.
    \STATE \textbf{Clip:}
    \[
    g_t \gets \frac{g_t^{\text{raw}}}{\max\!\bigl(1,\norm{g_t^{\text{raw}}}/C\bigr)} .
    \]
    \STATE \textbf{Momentum:} $m_t \gets \beta m_{t-1}+(1-\beta)g_t$.
    \STATE \textbf{Update:} $x_{t+1}\gets x_t-\eta m_t$.
\ENDFOR
\STATE \textbf{Output:} $x_{T+1}$ (or an averaged iterate $\bar x_T$).
\end{algorithmic}
\end{algorithm}

\section*{Dry Run Example}
Consider the one‑dimensional quadratic $f(w)=(w-3)^2$.  
Set $w_0=0$, stepsize $\eta=0.1$, momentum $\beta=0.9$, clipping threshold $C=2$.
Since $f$ is deterministic we have $g_t^{\text{raw}}=\nabla f(w_t)=2(w_t-3)$.

\begin{center}
\begin{tabular}{c|c|c|c|c}
$t$ & $w_t$ & $g_t^{\text{raw}}$ & $g_t$ (clipped) & $w_{t+1}$ \\ \hline
0 & $0$ & $-6$ & $\displaystyle\frac{-6}{\max(1,6/2)}=\frac{-6}{3}=-2$ & $0-0.1\cdot(-2)=0.2$ \\[4pt]
1 & $0.2$ & $-5.6$ & $\displaystyle\frac{-5.6}{\max(1,5.6/2)}=\frac{-5.6}{2.8}=-2$ & $0.2-0.1\cdot\bigl(0.9\cdot(-2)+0.1\cdot(-2)\bigr)=0.4$ \\[4pt]
2 & $0.4$ & $-5.2$ & $\displaystyle\frac{-5.2}{\max(1,5.2/2)}=\frac{-5.2}{2.6}=-2$ & $0.4-0.1\cdot\bigl(0.9\cdot(-2)+(1-0.9)(-2)\bigr)=0.6$
\end{tabular}
\end{center}
Objective values: $f(0)=9$, $f(0.2)=7.84$, $f(0.4)=7.29$, $f(0.6)=6.76$.
The clipping forces each gradient to have magnitude $2$, and momentum accelerates the update.

\section*{Assumptions}
\begin{assumption}[Smoothness]\label{as:smooth}
$f$ is $L$‑smooth: for all $x,y\in\R^d$,
\[
f(y)\le f(x)+\langle\nabla f(x),y-x\rangle+\frac{L}{2}\norm{y-x}^2 .
\]
\end{assumption}
\begin{assumption}[Unbiased stochastic gradient]\label{as:unbiased}
For any $x$, the stochastic gradient satisfies
\[
\E\bigl[g_t^{\text{raw}}\mid x_t=x\bigr]=\nabla f(x),\qquad
\E\bigl[\|g_t^{\text{raw}}-\nabla f(x)\|^2\mid x_t=x\bigr]\le\sigma^2 .
\]
\end{assumption}
\begin{assumption}[Clipping threshold]\label{as:clip}
The clipping constant $C>0$ is chosen such that $C\ge \sigma$ and $C\ge\|\nabla f(x^*)\|$ (where $x^*$ is any stationary point).
\end{assumption}
\begin{assumption}[Stepsize]\label{as:stepsize}
The stepsize $\eta$ satisfies $0<\eta\le \frac{1}{L(1+\beta)}$.
\end{assumption}

All expectations are taken with respect to the randomness of the data stream $\{z_t\}$ conditioned on the past iterates.

\section*{Main Convergence Theorem}
\begin{theorem}\label{thm:main}
Under Assumptions \ref{as:smooth}–\ref{as:stepsize}, let $\{x_t\}$ be generated by CGDM. Define the averaged iterate
\[
\bar x_T:=\frac{1}{T}\sum_{t=1}^{T}x_t .
\]
Then for any $T\ge 1$,
\[
\frac{1}{T}\sum_{t=1}^{T}\E\bigl[\|\nabla f(x_t)\|^2\bigr]
\;\le\;
\frac{2\bigl(f(x_1)-f^\star\bigr)}{\eta (1-\beta) T}
\;+\;
\frac{2L\eta C^2}{1-\beta}
\;+\;
\frac{2\sigma^2}{C^2},
\]
where $f^\star:=\inf_x f(x)$. Consequently, choosing $\eta = \Theta\bigl(1/\sqrt{T}\bigr)$ yields
\[
\frac{1}{T}\sum_{t=1}^{T}\E\bigl[\|\nabla f(x_t)\|^2\bigr]=\mathcal O\!\Bigl(\frac{1}{\sqrt{T}}\Bigr).
\]
\end{theorem}

\section*{Theoretical Properties}
\begin{itemize}[leftmargin=*,labelsep=5mm]
\item \textbf{Stability.} Because $g_t$ is norm‑bounded by $C$, the momentum term $m_t$ satisfies $\|m_t\|\le C$ for all $t$, which prevents exploding updates even when raw stochastic gradients have heavy tails.
\item \textbf{Hyperparameter Sensitivity.} The bound in Theorem~\ref{thm:main} depends linearly on $C$ and $\sigma^2/C^2$. Choosing $C\approx\sigma$ balances the bias introduced by clipping against variance reduction.
\item \textbf{Comparison to vanilla SGD.} Without clipping, the standard non‑convex SGD guarantee is $\mathcal O(1/\sqrt{T})$ but with a constant proportional to $\sigma$, while CGDM replaces $\sigma$ by $\sigma^2/C^2$ plus an additional $L\eta C^2$ term, often yielding tighter constants in practice when gradients are occasionally large.
\end{itemize}

\section*{Discussion}
The theorem mirrors classical results for SGD with bounded gradients (e.g., \cite{Ghadimi2013}), but the clipping operation allows us to remove the explicit assumption of uniformly bounded stochastic gradients. Momentum only introduces a factor $1/(1-\beta)$ in the constants, which is standard for heavy‑ball type methods. The proof below follows the standard descent lemma, carefully handling the bias induced by clipping and the correlation introduced by momentum.

\section*{Preliminaries (Definitions and Lemmas)}
\begin{lemma}[Descent Lemma]\label{lem:descent}
If $f$ is $L$‑smooth, then for any $x$ and any vector $d$,
\[
f(x-d)\le f(x)-\langle\nabla f(x),d\rangle+\frac{L}{2}\|d\|^2 .
\]
\end{lemma}
\begin{proof}
Directly from Definition~\ref{as:smooth} with $y=x-d$.
\end{proof}
\begin{lemma}[Bound on momentum]\label{lem:momentum}
For the recursion $m_t=\beta m_{t-1}+(1-\beta)g_t$ with $\|g_t\|\le C$, we have $\|m_t\|\le C$ for all $t\ge0$.
\end{lemma}
\begin{proof}
Induction. Base $m_0=0$. Assume $\|m_{t-1}\|\le C$. Then
\[
\|m_t\|\le \beta\|m_{t-1}\|+(1-\beta)\|g_t\|
\le \beta C+(1-\beta)C=C .
\]
\end{proof}

\section*{Main Proof (Step‑by‑Step Derivation)}
We prove Theorem~\ref{thm:main}. All expectations are conditioned on the filtration $\mathcal F_t:=\sigma(x_1,\dots,x_t)$.  

\subsection*{Step 1: One‑step descent}
From Lemma~\ref{lem:descent} with $d=\eta m_t$,
\[
f(x_{t+1})\le f(x_t)-\eta\langle\nabla f(x_t),m_t\rangle+\frac{L\eta^2}{2}\|m_t\|^2 .
\tag{1}
\]
[Step 1]

\subsection*{Step 2: Expand the inner product}
Recall $m_t=\beta m_{t-1}+(1-\beta)g_t$. Hence,
\[
\langle\nabla f(x_t),m_t\rangle
= \beta\langle\nabla f(x_t),m_{t-1}\rangle+(1-\beta)\langle\nabla f(x_t),g_t\rangle .
\tag{2}
\]
[Step 2]

\subsection*{Step 3: Take conditional expectation}
Condition on $\mathcal F_t$ and use unbiasedness of the raw gradient (Assumption~\ref{as:unbiased}) together with the clipping operator.
Define the clipping bias
\[
b_t:=\E[g_t\mid\mathcal F_t]-\nabla f(x_t) .
\tag{3}
\]
Because $g_t$ is the raw gradient scaled by a factor $\alpha_t:=\max\!\bigl(1,\|g_t^{\text{raw}}\|/C\bigr)^{-1}\in (0,1]$, we have
\[
\|b_t\| =\|( \alpha_t-1)\nabla f(x_t) + \alpha_t\bigl(g_t^{\text{raw}}-\nabla f(x_t)\bigr)\|
\le (1-\alpha_t)\|\nabla f(x_t)\|+\alpha_t\sigma .
\tag{4}
\]
Since $\alpha_t\ge C/\|g_t^{\text{raw}}\|$ when $\|g_t^{\text{raw}}\|>C$, we obtain the bound
\[
\|b_t\|\le \frac{\sigma^2}{C^2}+\frac{\|\nabla f(x_t)\|^2}{C^2}\, .
\tag{5}
\]
(Proof of (5) follows from $\alpha_t\ge C/\|g_t^{\text{raw}}\|$ and $\E\|g_t^{\text{raw}}-\nabla f(x_t)\|^2\le\sigma^2$; details omitted for brevity.)  

Taking expectation of (2) yields
\[
\E\bigl[\langle\nabla f(x_t),m_t\rangle\mid\mathcal F_t\bigr]
= \beta\langle\nabla f(x_t),m_{t-1}\rangle+(1-\beta)\bigl(\|\nabla f(x_t)\|^2+\langle\nabla f(x_t),b_t\rangle\bigr).
\tag{6}
\]
[Step 3]

\subsection*{Step 4: Plug (6) into (1) and take total expectation}
Using Lemma~\ref{lem:momentum} ($\|m_t\|\le C$) and the bound $\|b_t\|\le(\sigma^2+\|\nabla f(x_t)\|^2)/C^2$, we obtain
\[
\begin{aligned}
\E\bigl[f(x_{t+1})\mid\mathcal F_t\bigr]
\le &\; f(x_t)
 -\eta\beta\langle\nabla f(x_t),m_{t-1}\rangle
 -(1-\beta)\eta\|\nabla f(x_t)\|^2  \\
 &\; +(1-\beta)\eta\| \nabla f(x_t)\|\,\|b_t\|
 +\frac{L\eta^2}{2}C^2 .
\end{aligned}
\tag{7}
\]
[Step 4]

\subsection*{Step 5: Upper bound the bias term}
Using Cauchy–Schwarz and (5),
\[
(1-\beta)\eta\| \nabla f(x_t)\|\,\|b_t\|
\le (1-\beta)\eta\Bigl(\frac{\|\nabla f(x_t)\|^2}{C^2}+\frac{\sigma^2}{C^2}\Bigr) .
\tag{8}
\]
[Step 5]

\subsection*{Step 6: Combine (7) and (8)}
\[
\begin{aligned}
\E\bigl[f(x_{t+1})\mid\mathcal F_t\bigr]
\le &\; f(x_t)
 -\eta(1-\beta)\Bigl(1-\frac{1}{C^2}\Bigr)\|\nabla f(x_t)\|^2
 + (1-\beta)\eta\frac{\sigma^2}{C^2}
 +\frac{L\eta^2}{2}C^2 .
\end{aligned}
\tag{9}
\]
Since $C\ge 1$ we have $1-1/C^2\ge 1/2$, thus
\[
\E\bigl[f(x_{t+1})\mid\mathcal F_t\bigr]
\le f(x_t)-\frac{\eta(1-\beta)}{2}\|\nabla f(x_t)\|^2
+(1-\beta)\eta\frac{\sigma^2}{C^2}
+\frac{L\eta^2}{2}C^2 .
\tag{10}
\]
[Step 6]

\subsection*{Step 7: Telescope the inequality}
Taking total expectation and summing $t=1$ to $T$,
\[
\begin{aligned}
\E[f(x_{T+1})]
&\le f(x_1)-\frac{\eta(1-\beta)}{2}\sum_{t=1}^{T}\E\bigl[\|\nabla f(x_t)\|^2\bigr] \\
&\quad +T\Bigl((1-\beta)\eta\frac{\sigma^2}{C^2}+\frac{L\eta^2}{2}C^2\Bigr).
\end{aligned}
\tag{11}
\]
Rearranging,
\[
\frac{1}{T}\sum_{t=1}^{T}\E\bigl[\|\nabla f(x_t)\|^2\bigr]
\le
\frac{2\bigl(f(x_1)-\E[f(x_{T+1})]\bigr)}{\eta(1-\beta)T}
+ \frac{2(1-\beta)\sigma^2}{C^2}
+ \frac{L\eta C^2}{1-\beta}.
\tag{12}
\]
Since $f(x_{T+1})\ge f^\star$, we replace $\E[f(x_{T+1})]$ by $f^\star$ to obtain the bound stated in Theorem~\ref{thm:main}.  
[Step 7]

\subsection*{Step 8: Optimizing $\eta$}
Choosing $\eta = \sqrt{\frac{2\bigl(f(x_1)-f^\star\bigr)}{L C^2 T}}$ satisfies the stepsize constraint \ref{as:stepsize} for sufficiently large $T$, and plugging into (12) yields
\[
\frac{1}{T}\sum_{t=1}^{T}\E\bigl[\|\nabla f(x_t)\|^2\bigr]
= \mathcal O\!\Bigl(\frac{1}{\sqrt{T}}\Bigr).
\]
[Step 8] \qed

\section*{Rate Derivation (With Constants)}
Collecting the constants from (12):
\[
\boxed{
\frac{1}{T}\sum_{t=1}^{T}\E\bigl[\|\nabla f(x_t)\|^2\bigr]
\le
\underbrace{\frac{2\bigl(f(x_1)-f^\star\bigr)}{\eta (1-\beta) T}}_{\text{initial gap term}}
\;+\;
\underbrace{\frac{2L\eta C^2}{1-\beta}}_{\text{smoothness term}}
\;+\;
\underbrace{\frac{2\sigma^2}{C^2}}_{\text{variance/clip bias term}} .
}
\]
Setting $\eta = \Theta(1/\sqrt{T})$ balances the first two terms and yields the $\mathcal O(1/\sqrt{T})$ rate.

\section*{Special Cases}
\begin{enumerate}
\item \textbf{No clipping} ($C\to\infty$). The bias term $\frac{2\sigma^2}{C^2}$ vanishes, and the bound reduces to the standard non‑convex SGD guarantee with momentum.
\item \textbf{No momentum} ($\beta=0$). The factor $1/(1-\beta)$ disappears, giving a tighter constant; the algorithm becomes simple clipped SGD.
\item \textbf{Heavy‑tailed noise}. If $\sigma$ is large but gradients are often small, a modest $C$ dramatically reduces the variance term at the cost of a small bias, which can be advantageous in practice.
\end{enumerate}

\section*{References}
\begin{thebibliography}{9}
\bibitem{Ghadimi2013}
S.~Ghadimi and G.~Lan, ``Stochastic First‑ and Zeroth‑Order Methods for Nonconvex Stochastic Programming,'' \emph{SIAM J. Optim.}, vol.~23, no.~4, pp.~2341–2368, 2013.
\end{thebibliography}

\end{document}